% GENERATED FILE. Edit canonical lessons, then run npm run generate:books.
% canonical-source-hash: fnv1a64:6f3f458de8040910
% canonical-lessons: BN-C401-dokkhin, BN-C401-purbo, BN-C401-poshchim, BN-C401-kompas, BN-C401-digonto

\chapter{Points of the Compass}
\label{ch:bn-points-of-the-compass}
\hlchaptermodality{401}

\begin{chapteropening}
\textbf{By the end of this chapter:} \emph{I can say which way is south, east and west, read a compass and look to the horizon.}

Complete the last lesson of chapter 401: I can say which way is south, east and west, read a compass and look to the horizon.
\end{chapteropening}

\section[\texorpdfstring{\bn{দক্ষিণ} (dôkkhin)}{দক্ষিণ (dôkkhin)}]{\bn{দক্ষিণ} (dôkkhin) — south}

\label{lesson:BN-C401-dokkhin}



\begin{quote}
Before the new one: say the Bengali for an instruction, then the Bengali for to point (to).
\end{quote}

\subsection*{You'll want to know: \bn{দক্ষিণ}}
\textbf{\bn{দক্ষিণ}} — \emph{dôkkhin} — \textquotedblleft{}south\textquotedblright{}.

\begin{quote}
\bn{সমুদ্র} \bn{এই} \bn{গ্রামের} \bn{দক্ষিণে।}
\end{quote}

\begin{quote}
\emph{sômudrô ei grāmer dôkkhine.}
\end{quote}

\begin{quote}
The sea lies to the south of this village.
\end{quote}

\bn{ক্ষ} is said kkh in the middle of a word. \bn{উত্তর}, north, is also the word for an answer.

\begin{cousinweb}
\textbf{\bn{দক্ষিণ}} (\emph{dôkkhin}) is Sanskrit \emph{dakṣiṇa}, \textquotedblleft{}on the right hand\textquotedblright{} and so \textquotedblleft{}southern\textquotedblright{} (Monier-Williams): facing the sunrise, the south is on your right.
\end{cousinweb}

\begin{grammarlens}[title={What you've built}]
One more word for getting around.
\end{grammarlens}

\subsection*{Your turn}
\begin{itemize}
\raggedright
  \item \emph{Say it:} \emph{dôkkhin}
  \item \emph{Say it:} \emph{dôkkhin}, once more
  \item \emph{Say it:} \emph{sômudrô ei grāmer dôkkhine.} — The sea lies to the south of this village.
  \item \emph{Recall:} say the Bengali for an instruction, then the Bengali for to point (to), then say \emph{dôkkhin} again
\end{itemize}

\begin{tcolorbox}[breakable,colback=teal!4,colframe=teal!35!black,title={Before you move on}]
What does \bn{দক্ষিণ} mean? (\textquotedblleft{}South\textquotedblright{}.) Say it once more.
\end{tcolorbox}
\section[\texorpdfstring{\bn{পূর্ব} (pūrbô)}{পূর্ব (pūrbô)}]{\bn{পূর্ব} (pūrbô) — east}

\label{lesson:BN-C401-purbo}



\begin{quote}
Before the new one: say the Bengali for to point (to), then the Bengali for south.
\end{quote}

\subsection*{You'll want to know: \bn{পূর্ব}}
\textbf{\bn{পূর্ব}} — \emph{pūrbô} — \textquotedblleft{}east\textquotedblright{}.

\begin{quote}
\bn{ভোরে} \bn{পূর্ব} \bn{আকাশ} \bn{লাল} \bn{হল।}
\end{quote}

\begin{quote}
\emph{bhore pūrbô ākāsh lāl hôlô.}
\end{quote}

\begin{quote}
At dawn the eastern sky turned red.
\end{quote}

Also former, earlier, in careful speech: \bn{পূর্বে}, before.

\begin{cousinweb}
\textbf{\bn{পূর্ব}} (\emph{pūrbô}) is Sanskrit \emph{pūrva}, \textquotedblleft{}being before, in front\textquotedblright{}, and so \textquotedblleft{}eastern\textquotedblright{} and \textquotedblleft{}earlier\textquotedblright{} (Monier-Williams): the east was the side you faced.
\end{cousinweb}

\begin{grammarlens}[title={What you've built}]
One more word for getting around.
\end{grammarlens}

\subsection*{Your turn}
\begin{itemize}
\raggedright
  \item \emph{Say it:} \emph{pūrbô}
  \item \emph{Say it:} \emph{pūrbô}, once more
  \item \emph{Say it:} \emph{bhore pūrbô ākāsh lāl hôlô.} — At dawn the eastern sky turned red.
  \item \emph{Recall:} say the Bengali for to point (to), then the Bengali for south, then say \emph{pūrbô} again
\end{itemize}

\begin{tcolorbox}[breakable,colback=teal!4,colframe=teal!35!black,title={Before you move on}]
What does \bn{পূর্ব} mean? (\textquotedblleft{}East\textquotedblright{}.) Say it once more.
\end{tcolorbox}
\section[\texorpdfstring{\bn{পশ্চিম} (pôshchim)}{পশ্চিম (pôshchim)}]{\bn{পশ্চিম} (pôshchim) — west}

\label{lesson:BN-C401-poshchim}



\begin{quote}
Before the new one: say the Bengali for south, then the Bengali for east.
\end{quote}

\subsection*{You'll want to know: \bn{পশ্চিম}}
\textbf{\bn{পশ্চিম}} — \emph{pôshchim} — \textquotedblleft{}west\textquotedblright{}.

\begin{quote}
\bn{সন্ধ্যায়} \bn{আমরা} \bn{পশ্চিম} \bn{দিকে} \bn{হাঁটলাম।}
\end{quote}

\begin{quote}
\emph{sôndhyāy āmrā pôshchim dike hā̃ṭlām.}
\end{quote}

\begin{quote}
In the evening we walked westward.
\end{quote}

\begin{cousinweb}
\textbf{\bn{পশ্চিম}} (\emph{pôshchim}) is Sanskrit \emph{paścima}, \textquotedblleft{}being behind, later\textquotedblright{}, and so \textquotedblleft{}western\textquotedblright{} (Monier-Williams): with the sunrise in front, the west lies behind you.
\end{cousinweb}

\begin{grammarlens}[title={What you've built}]
One more word for getting around.
\end{grammarlens}

\subsection*{Your turn}
\begin{itemize}
\raggedright
  \item \emph{Say it:} \emph{pôshchim}
  \item \emph{Say it:} \emph{pôshchim}, once more
  \item \emph{Say it:} \emph{sôndhyāy āmrā pôshchim dike hā̃ṭlām.} — In the evening we walked westward.
  \item \emph{Recall:} say the Bengali for south, then the Bengali for east, then say \emph{pôshchim} again
\end{itemize}

\begin{tcolorbox}[breakable,colback=teal!4,colframe=teal!35!black,title={Before you move on}]
What does \bn{পশ্চিম} mean? (\textquotedblleft{}West\textquotedblright{}.) Say it once more.
\end{tcolorbox}
\section[\texorpdfstring{\bn{কম্পাস} (kômpās)}{কম্পাস (kômpās)}]{\bn{কম্পাস} (kômpās) — a compass}

\label{lesson:BN-C401-kompas}



\begin{quote}
Before the new one: say the Bengali for east, then the Bengali for west.
\end{quote}

\subsection*{You'll want to know: \bn{কম্পাস}}
\textbf{\bn{কম্পাস}} — \emph{kômpās} — \textquotedblleft{}a compass\textquotedblright{}.

\begin{quote}
\bn{বনে} \bn{আমরা} \bn{কম্পাস} \bn{দেখে} \bn{পথ} \bn{খুঁজলাম।}
\end{quote}

\begin{quote}
\emph{bône āmrā kômpās dekhe pôth khũjlām.}
\end{quote}

\begin{quote}
In the forest we looked for the path by the compass.
\end{quote}

\begin{cousinweb}
\textbf{\bn{কম্পাস}} (\emph{kômpās}) is borrowed from English \emph{compass}.
\end{cousinweb}

\begin{grammarlens}[title={What you've built}]
One more word for getting around.
\end{grammarlens}

\subsection*{Your turn}
\begin{itemize}
\raggedright
  \item \emph{Say it:} \emph{kômpās}
  \item \emph{Say it:} \emph{kômpās}, once more
  \item \emph{Say it:} \emph{bône āmrā kômpās dekhe pôth khũjlām.} — In the forest we looked for the path by the compass.
  \item \emph{Recall:} say the Bengali for east, then the Bengali for west, then say \emph{kômpās} again
\end{itemize}

\begin{tcolorbox}[breakable,colback=teal!4,colframe=teal!35!black,title={Before you move on}]
What does \bn{কম্পাস} mean? (\textquotedblleft{}A compass\textquotedblright{}.) Say it once more.
\end{tcolorbox}
\section[\texorpdfstring{\bn{দিগন্ত} (digôntô)}{দিগন্ত (digôntô)}]{\bn{দিগন্ত} (digôntô) — the horizon}

\label{lesson:BN-C401-digonto}



\begin{quote}
Before the new one: say the Bengali for west, then the Bengali for a compass.
\end{quote}

\subsection*{You'll want to know: \bn{দিগন্ত}}
\textbf{\bn{দিগন্ত}} — \emph{digôntô} — \textquotedblleft{}the horizon\textquotedblright{}.

\begin{quote}
\bn{সৈকত} \bn{থেকে} \bn{আমরা} \bn{দিগন্তে} \bn{একটা} \bn{জাহাজ} \bn{দেখলাম।}
\end{quote}

\begin{quote}
\emph{soikôt theke āmrā digônte ekṭā jāhāj dekhlām.}
\end{quote}

\begin{quote}
From the beach we saw a ship on the horizon.
\end{quote}

\begin{cousinweb}
\textbf{\bn{দিগন্ত}} (\emph{digôntô}) is Sanskrit \emph{dig-anta}, \textquotedblleft{}the end of the horizon, a far distance\textquotedblright{} (Monier-Williams), from \emph{diś}, a quarter of the sky, and \emph{anta}, an end.
\end{cousinweb}

\begin{grammarlens}[title={What you've built}]
One more word for getting around.
\end{grammarlens}

\subsection*{Your turn}
\begin{itemize}
\raggedright
  \item \emph{Say it:} \emph{digôntô}
  \item \emph{Say it:} \emph{digôntô}, once more
  \item \emph{Say it:} \emph{soikôt theke āmrā digônte ekṭā jāhāj dekhlām.} — From the beach we saw a ship on the horizon.
  \item \emph{Recall:} say the Bengali for west, then the Bengali for a compass, then say \emph{digôntô} again
\end{itemize}

\begin{tcolorbox}[breakable,colback=teal!4,colframe=teal!35!black,title={Before you move on}]
What does \bn{দিগন্ত} mean? (\textquotedblleft{}The horizon\textquotedblright{}.) Say it once more.
\end{tcolorbox}
